\section{干预思维的证据车道}

本章明确使用本书的证据政策。

\begin{itemize}[nosep]
  \item \textbf{强证据车道。} 来自意识与神经动力学研究的状态转变证据，支持这样一个一般性主张：系统能够占据不同的全局机制，并通过结构化转变在其间移动。这使我们有理由把干预视为一种机制塑形，而不只是自我对话。
  \item \textbf{中等证据车道。} 亚稳态、吸引子、积分器、因果涌现与非线性控制等视角，支持这样一种设计观念：反复的局部输入能够稳定或重定向未来行为。
  \item \textbf{推测性车道。} 前沿意识物理学方案也许能激发关于大尺度组织的假设，但它们并不足以为本章中的实际干预规则提供正当性。
  \item \textbf{设计推断。} 下文给出的具体协议属于应用综合。它们受强证据车道与中等证据车道启发，但仍然是设计构造，而不是直接的实验定理。
\end{itemize}

\begin{proposition}[干预应瞄准转变几何结构]
当一种失败反复出现时，首要的设计问题通常应当是：哪一种改变能够以最低成本，在真实的决策窗口中提高进入目标状态的概率？
\end{proposition}

\begin{discussion}
这一命题比“让任务更容易”更强。它要求我们用转变语言来评估干预。如果当前状态是一个局部稳定的盆地，那么有效干预要么降低障碍，要么改变尝试移动的时机，要么改变局部环境，使旧状态变得不再那么容易自我恢复。这一逻辑既与本书的状态转变框架一致，也与更广泛的动力系统语言一致，包括吸引子、亚稳态以及类似分岔的机制切换。
\end{discussion}